% =====================================================================
% frontmatter/resumo.tex
% PT: Resumo em português europeu. Tradução fiel do Abstract, também num
%     único parágrafo com 150 a 300 palavras (confirme o limite no
%     regulamento em vigor). As palavras-chave (normalmente 3 a 6) vão no
%     fim, na linha "Palavras-chave", separadas por ponto e vírgula, pela
%     mesma ordem das Keywords.
% EN: Abstract in European Portuguese. A faithful translation of the
%     Abstract, also a single paragraph of 150 to 300 words (check the
%     limit in the current regulations). The keywords (usually 3 to 6) go
%     at the end, on the "Palavras-chave" line, separated by semicolons, in
%     the same order as the English keywords.
% =====================================================================

% PT: O ambiente otherlanguage aplica a hifenização portuguesa.
% EN: The otherlanguage environment applies Portuguese hyphenation.
\begin{otherlanguage}{portuguese}
\chapter*{Resumo}
\addcontentsline{toc}{chapter}{Resumo}

% PT: Texto de exemplo em português europeu: substitua pelo seu resumo.
% EN: Example text in European Portuguese: replace it with your abstract.
Este é um resumo de exemplo. Numa dissertação real, este parágrafo
sintetiza todo o trabalho para um leitor que não vai ler mais nada. Começa
com uma ou duas frases de contexto que explicam a relevância do tema e
indica depois o problema ou a lacuna a que a dissertação dá resposta. Em
seguida, descreve a abordagem em termos concretos: o método, os dados ou o
caso de estudo e as ferramentas utilizadas. Apresenta depois os principais
resultados, de preferência com os valores mais importantes, e termina com a
conclusão principal e o seu significado para a prática ou para investigação
futura. O resumo não deve conter citações, figuras, tabelas nem siglas por
definir, e todas as afirmações devem estar fundamentadas no texto
principal. Escreva-o no fim, quando os resultados e as conclusões estiverem
fechados, e respeite o limite de palavras definido pela escola.

\vspace{1em}
% PT: Palavras-chave em português (3 a 6), separadas por ponto e vírgula.
% EN: Keywords in Portuguese (3 to 6), separated by semicolons.
\noindent\textbf{Palavras-chave:} Palavra-chave Um; Palavra-chave Dois; Palavra-chave Três; Palavra-chave Quatro
\end{otherlanguage}
